\chapter{电子效应与反应活性中间体}\label{ch:b1:electronic-effects}

三氟乙酸 \ce{CF3COOH} 的酸性约为乙酸
\ce{CH3COOH} 的一万倍，尽管其酸性的
O--H 键与氟原子之间相隔三个键。苯酚的酸性约为乙醇的
一百万倍，而一个与苯环相连的氨分子（即苯胺）
的碱性仅为与甲基相连的氨分子的百万分之一。有机化学中
此类对比屡见不鲜，其中大部分可由两种效应解释：
一是\omterm{def:b1:periodicity:electronegativity}{电负性}原子沿
$\sigma$ 键的拉电子作用，二是电子对在 $\pi$
体系中离域。同样这两种效应决定哪些键断裂、哪些中间体
生成以及试剂在何处进攻，它们是后续各章
\omterm{def:b1:elementary-steps:elementary-step}{反应机理}的“语法”。

\begin{recall}
\omterm{def:b1:periodicity:electronegativity}{电负性}及其变化趋势（见\cref{ch:b1:periodicity}）；\omterm{def:b1:lewis-resonance-vsepr:lewis}{路易斯
结构}、\omterm{def:b1:lewis-resonance-vsepr:formal-charge}{形式电荷}与共振（见\cref{ch:b1:lewis-resonance-vsepr}）；
$\mathrm{p}K_a$ 与酸、碱的强弱
（见\cref{ch:b1:predominant-reaction}）。第 1 册（十二年级）曾
画出从\omterm{def:b1:electronic-effects:nucleophile}{亲核试剂}指向\omterm{def:b1:electronic-effects:nucleophile}{亲电试剂}的\omterm{def:b1:electronic-effects:curly-arrow}{弯箭头}。
\end{recall}

\section{诱导效应}

\begin{definition}[诱导效应]\label{def:b1:electronic-effects:inductive}
\emph{诱导效应}\index{诱导效应}是电子云密度沿
$\sigma$ 键向\omterm{def:b1:periodicity:electronegativity}{电负性}较强的原子或
基团偏移。\omterm{def:b1:electronic-effects:mesomeric}{吸电子基团}（F、Cl、O、\ce{NO2}、\ce{CF3}）具有 $-I$
效应；烷基及带负电的原子推电子，具有 $+I$
效应。该效应随中间键数的增多而迅速减弱。
\end{definition}

\begin{proposition}[诱导效应与酸性]\label{prop:b1:electronic-effects:inductive-pka}
羧酸酸性位点附近的\omterm{def:b1:electronic-effects:mesomeric}{吸电子基团}
可通过分散羧酸根离子的负电荷使其稳定，
并降低 $\mathrm{p}K_a$：\omterm{def:b1:electronic-effects:mesomeric}{吸电子基团}越多、越近，
酸性就越强。
\end{proposition}

\begin{proof}
$\mathrm{p}K_a$ 表示下列平衡的位置：$\mathrm{RCOOH} \rightleftharpoons
\mathrm{RCOO^-} + \ce{H+}$。能将电子云从
羧酸根处拉开的基团使其负电荷分散到更多原子上，
其能量降低的幅度大于中性酸；平衡向右移动，
$K_a$ 增大。实测数据证实了这一点。
\end{proof}

\begin{omfigure}
\begin{tikzpicture}
\begin{axis}[
  xbar, width=0.78\linewidth, height=5.2cm, xmin=0, xmax=5.4,
  symbolic y coords={TFA,TCA,DCA,MCA,AA}, ytick=data,
  yticklabels={\ce{CF3COOH}, \ce{CCl3COOH}, \ce{CHCl2COOH}, \ce{CH2ClCOOH}, \ce{CH3COOH}},
  xlabel={$\mathrm{p}K_a$（\qty{25}{\celsius}）}, bar width=9pt,
  nodes near coords, nodes near coords align={horizontal},
  every node near coord/.append style={font=\scriptsize},
  tick label style={font=\small}, label style={font=\small}, enlarge y limits=0.15]
\addplot[fill=omDef!45, draw=omDef] coordinates {(0.52,TFA) (0.51,TCA) (1.35,DCA) (2.87,MCA) (4.76,AA)};
\end{axis}
\end{tikzpicture}

{\small 卤代乙酸。每个氯原子都会
使 $\mathrm{p}K_a$ 降低；含三个氯原子时，
该酸的酸性约为乙酸的一万倍。三氟乙酸与三氯乙酸
的强度在此类\omterm{def:b1:predominant-reaction:strong-acid}{强酸}的测量
不确定度范围内相当。}
% ledger: pka:ethanoic, pka:chloroethanoic, pka:dichloroethanoic, pka:trichloroethanoic, pka:trifluoroethanoic
\end{omfigure}

烷基作用方向相反，影响较小：丙酸（$\mathrm{p}K_a
= 4.89$）的酸性略弱于乙酸（4.76），而甲酸（
3.74，不含烷基）的酸性强于上述两者。
% ledger: pka:propanoic, pka:methanoic

\section{共轭效应}

当带孤对电子或电荷的原子与 $\pi$
体系相连时，其电子可与该体系共享。\omterm{def:b1:lewis-resonance-vsepr:resonance}{共振
结构}（见\cref{ch:b1:lewis-resonance-vsepr}）描述这种共享。

\begin{definition}[共轭效应、给电子基团与吸电子基团]\label{def:b1:electronic-effects:mesomeric}
\emph{共轭效应}\index{共轭效应}是电子云密度
沿共轭 $\pi$ 体系的偏移，以\omterm{def:b1:lewis-resonance-vsepr:resonance}{共振
结构}表示。\emph{给电子基团}\index{给电子基团}向体系提供一对孤对电子，
具有 $+M$ 效应（\ce{-OH}、\ce{-O^-}、\ce{-NH2}、\ce{-OR}）；一个
\emph{吸电子基团}\index{吸电子基团}通过
$\pi$ 键从体系中将电子拉向\omterm{def:b1:periodicity:electronegativity}{电负性}原子，具有 $-M$ 效应（\ce{-NO2}、
\ce{-C=O}、\ce{-C#N}）。
\end{definition}

\begin{omfigure}
\setchemfig{atom sep=2em}
\schemestart
\chemfig{CH_3-C(=[1]O)-[7]O^{-}}
\arrow{<->}
\chemfig{CH_3-C(-[1]O^{-})=[7]O}
\schemestop
\hspace{2.5em}
\schemestart
\chemfig{*6(=-=(-O^{-})-=-)}
\arrow{<->}
\chemfig{*6((=O)-=-\chemabove{}{\scriptstyle\ominus}-=-)}
\schemestop

\medskip
\schemestart
\chemfig{H_3C-C(=[1]O)-[7]NH_2}
\arrow{<->}
\chemfig{H_3C-C(-[1]O^{-})=[7]NH_2^{+}}
\schemestop

{\small 乙酸根离子的\omterm{def:b1:lewis-resonance-vsepr:resonance}{共振结构}（电荷由两个等价氧原子共享）、
苯氧负离子的\omterm{def:b1:lewis-resonance-vsepr:resonance}{共振结构}（电荷分散到环上；所示三种环结构之一）以及
酰胺的\omterm{def:b1:lewis-resonance-vsepr:resonance}{共振结构}（氮原子的孤对电子与 C=O 共享，
使酰胺成为\omterm{def:b1:predominant-reaction:strong-acid}{弱碱}）。}
\end{omfigure}

\begin{proposition}[共轭效应与酸性]\label{prop:b1:electronic-effects:mesomeric-pka}
共轭碱能靠共振分散电荷的酸，
比共轭碱不能分散电荷的酸更强：羧酸（$\mathrm{p}K_a$ 约为 4--5）
远强于醇（约 16），苯酚（10.0）介于两者之间。
\omterm{def:b1:electronic-effects:mesomeric}{吸电子基团}位于苯酚 OH 的邻位或对位时，
对其酸性的增强作用比位于间位时更大。
\end{proposition}

\begin{proof}
在醇盐中，电荷位于一个氧原子上；在羧酸根中，电荷由两个氧原子共享；
在苯氧负离子中，电荷由氧原子与三个环碳原子共享，
环碳的\omterm{def:b1:periodicity:electronegativity}{电负性}小于氧，稳定化程度较小。
位于邻位或对位碳上的 \ce{NO2} 基团，
可通过与 \ce{C=N+(-O^-)-O^-} 的\omterm{def:b1:lewis-resonance-vsepr:resonance}{共振结构}自身承担电荷；在间位上，
没有任何\omterm{def:b1:lewis-resonance-vsepr:resonance}{共振结构}能使电荷转移到带 \ce{NO2} 的碳上，
只剩\omterm{def:b1:electronic-effects:inductive}{诱导效应}起作用。测得的 pKa 值为
2-硝基苯酚 7.23、4-硝基苯酚 7.16、3-硝基苯酚 8.36，
而苯酚为 9.99，乙醇为 15.9。
\end{proof}
% ledger: pka:ethanol, pka:phenol, pka:2-nitrophenol, pka:3-nitrophenol, pka:4-nitrophenol

同样的推理也适用于碱。甲胺（其铵离子的 $\mathrm{p}K_a$
为 10.66）比氨（9.25，见\cref{ch:b1:predominant-reaction}）是更强的碱，
甲基将电子推向氮原子；在苯胺中，孤对电子
与苯环共享，苯胺的铵离子酸性要强得多（4.60）。
吡啶（5.23）保留了其孤对电子，但该氮原子处于平面环中，
具有更多的 \textit{s} 成分，因此其孤对电子的可利用性低于胺中的氮。
% ledger: pka:methylammonium, pka:anilinium, pka:pyridinium, dfg:NH4+_ao

\section{断键：活性中间体}

\begin{definition}[均裂与异裂]\label{def:b1:electronic-effects:cleavage}
\omterm{def:b1:lewis-resonance-vsepr:lewis}{共价键}断裂时，若每个原子各保留一个成键电子，则发生\emph{均裂}\index{均裂}，
生成两个\omterm{def:b1:electronic-effects:intermediates}{自由基}；
若一个原子保留两个电子，则发生\emph{异裂}\index{异裂}，
生成一个阳离子和一个阴离子。
\end{definition}

\begin{definition}[活性中间体]\label{def:b1:electronic-effects:intermediates}
称在机理的某一步中生成、又在之后某一步中被消耗的物种为
\emph{活性中间体}\index{活性中间体}，其活性太高，难以积累
（见\cref{ch:b1:elementary-steps}）。
\emph{碳正离子}\index{碳正离子}是具有三个键、带正电荷的碳原子
（\omterm{def:b1:quantum-numbers:configuration}{价电子}数为六个）；
\emph{碳负离子}\index{碳负离子}是具有三个键、一对孤对电子
并带负电荷的碳原子；\emph{自由基}\index{自由基}是含有
一个\omterm{def:b1:quantum-numbers:unpaired}{未成对电子}的原子。
\end{definition}

\begin{definition}[碳原子的类别]\label{def:b1:electronic-effects:carbon-class}
与一个其他碳原子相连的碳称为
\emph{伯碳}\index{伯碳}，与两个碳原子相连的称为
\emph{仲碳}\index{仲碳}，与三个碳原子相连的称为
\emph{叔碳}\index{叔碳}。\omterm{def:b1:electronic-effects:intermediates}{碳正离子}、\omterm{def:b1:electronic-effects:intermediates}{自由基}或卤代烷
的类别取决于所涉及的碳原子。
\end{definition}

\begin{omfigure}
\begin{tikzpicture}[font=\small]
  % carbocation: planar, empty p orbital
  \begin{scope}
    \fill[omProp!15, draw=omProp] (0,0.15) ellipse (0.18 and 0.75);
    \fill[omProp!15, draw=omProp] (0,-0.15) ellipse (0.18 and 0.75);
    \draw[thick] (0,0) -- (-1.0,-0.25) node[left] {R};
    \draw[thick] (0,0) -- (0.95,-0.35) node[right] {R};
    \draw[thick] (0,0) -- (0.35,0.45) node[above right] {R};
    \node[fill=white, inner sep=1pt] at (0,0) {\ce{C+}};
    \node at (0,-1.75) {碳正离子：平面形，};
    \node at (0,-2.15) {空的 \textit{p} 轨道};
  \end{scope}
  % radical: nearly planar, one electron
  \begin{scope}[xshift=4.6cm]
    \draw[omMeth] (0,0.15) ellipse (0.18 and 0.75);
    \draw[omMeth] (0,-0.15) ellipse (0.18 and 0.75);
    \fill (0,0.55) circle (0.05);
    \draw[thick] (0,0) -- (-1.0,-0.25) node[left] {R};
    \draw[thick] (0,0) -- (0.95,-0.35) node[right] {R};
    \draw[thick] (0,0) -- (0.35,0.45) node[above right] {R};
    \node[fill=white, inner sep=1pt] at (0,0) {C};
    \node at (0,-1.75) {自由基：近似平面形，};
    \node at (0,-2.15) {一个电子};
  \end{scope}
  % carbanion: pyramidal with a lone pair
  \begin{scope}[xshift=9.2cm]
    \draw[thick] (0,0) -- (-0.95,-0.55) node[left] {R};
    \draw[thick] (0,0) -- (0.95,-0.55) node[right] {R};
    \draw[thick] (0,0) -- (0.2,-0.9) node[below] {R};
    \fill[omDef!15, draw=omDef] (0,0.55) ellipse (0.22 and 0.45);
    \fill (-0.07,0.65) circle (0.045); \fill (0.07,0.65) circle (0.045);
    \node[fill=white, inner sep=1pt] at (0,0) {\ce{C-}};
    \node at (0,-1.75) {碳负离子：棱锥形，};
    \node at (0,-2.15) {孤对电子};
  \end{scope}
\end{tikzpicture}

{\small 碳的三种\omterm{def:b1:electronic-effects:intermediates}{活性中间体}的几何构型（R：H 或
烷基）。平面形的\omterm{def:b1:electronic-effects:intermediates}{碳正离子}两面均可被进攻，
这一点对
在\cref{ch:b1:nucleophilic-substitution}中的立体化学至关重要。}
\end{omfigure}

\begin{proposition}[碳正离子的稳定性]\label{prop:b1:electronic-effects:carbocation-order}
\omterm{def:b1:electronic-effects:mesomeric}{给电子基团}可使\omterm{def:b1:electronic-effects:intermediates}{碳正离子}稳定，其稳定性顺序为叔
$>$ 仲 $>$ 伯 $>$ 甲基；与双键相邻的\omterm{def:b1:electronic-effects:intermediates}{碳正离子}
（烯丙基型）或与苯环相连的（苄基型）\omterm{def:b1:electronic-effects:intermediates}{碳正离子}，还可通过
共振得到进一步稳定，与带孤对电子的原子相邻者（\ce{R-O-CH2+}）
则更为稳定。
\end{proposition}

\begin{proof}
每个烷基都将电子密度推向\omterm{def:b1:lewis-resonance-vsepr:hypervalent}{缺电子的}
碳（$+I$ 效应，以及其 C--H 键与
空轨道之间的相关相互作用，见本科二年级分册）。烯丙基型\omterm{def:b1:electronic-effects:intermediates}{碳正离子}有
两种\omterm{def:b1:lewis-resonance-vsepr:resonance}{共振结构}，电荷分别位于两端的碳上；苄基型\omterm{def:b1:electronic-effects:intermediates}{碳正离子}有四种，
其中三种位于环上；氧的孤对电子与\omterm{def:b1:electronic-effects:intermediates}{碳正离子}形成第四根键，
\ce{R-O+=CH2}，
其中每个原子都满足八隅体。
\end{proof}

\omterm{def:b1:electronic-effects:intermediates}{碳负离子}的顺序则相反（甲基 $>$ 伯 $>$ 仲
$>$ 叔），并由\omterm{def:b1:electronic-effects:mesomeric}{吸电子基团}稳定；\omterm{def:b1:electronic-effects:intermediates}{自由基}的顺序与
\omterm{def:b1:electronic-effects:intermediates}{碳正离子}相同，但影响较弱。

\section{弯箭头、亲核试剂和亲电试剂}

\begin{definition}[弯箭头]\label{def:b1:electronic-effects:curly-arrow}
在\omterm{def:b1:elementary-steps:elementary-step}{反应机理}中，\emph{弯箭头}\index{弯箭头}表示
电子对的移动：它从孤对电子或化学键出发，止于
一个原子（在此处形成孤对电子或新键）或两个原子之间（形成新键）。
\emph{单钩箭头}\index{单钩箭头}表示单个电子的移动，
用于\omterm{def:b1:electronic-effects:intermediates}{自由基}反应。
\end{definition}

\begin{method}[画弯箭头]\label{met:b1:electronic-effects:curly-arrows}
\begin{enumerate}
  \item 从电子出发：孤对电子、$\pi$ 键或 $\sigma$
        键，绝不从正电荷或不含电子的原子出发。
  \item 止于电子对的去处：在两原子之间形成新键，或在键断裂时止于
        \omterm{def:b1:periodicity:electronegativity}{电负性}较大的原子上的孤对电子。
  \item 每一步后计算电荷：向新键提供孤对电子的原子获得 $+1$，
        以孤对形式接受电子对的原子获得 $-1$；
        总电荷守恒。
  \item C、N、O 上的电子数绝不能超过八隅体：电子对到达时，
        必须有另一对离开。
\end{enumerate}
\end{method}

\begin{omfigure}
\setchemfig{atom sep=2.2em}
\schemestart
\chemfig{H_3C-C(-[2]H)(-[6]H)-[@{b}]@{br}Br}
\arrow{->[异裂]}[,1.5]
\chemfig{H_3C-C(-[2]H)(-[6]H)^{+}}\+\chemfig{Br^{-}}
\schemestop
\chemmove{\draw[->, shorten <=2pt, shorten >=2pt](b) .. controls +(90:6mm) and +(90:6mm) .. (br);}

\medskip
\schemestart
\chemfig{H@{o}O^{-}}\+\chemfig{@{h}H-@{oc}O-H}
\arrow{<=>}
\chemfig{H_2O}\+\chemfig{HO^{-}}
\schemestop
\chemmove{\draw[->, thick, shorten <=2pt, shorten >=2pt](o) .. controls +(-80:8mm) and +(-100:8mm) .. (h);}
\par\vspace{5mm}

{\small \omterm{def:b1:electronic-effects:curly-arrow}{弯箭头}。上图：C--Br 键的\omterm{def:b1:electronic-effects:cleavage}{异裂}，电子对
归于溴。下图：质子从水转移到氢氧根，写作
碱的孤对电子进攻氢原子（水中的 O--H 键
随后断裂，其电子对留在氧上；第二个箭头
留给\cref{exo:b1:electronic-effects:4}）。}
\end{omfigure}

\begin{definition}[亲核试剂、亲电试剂、离去基团]\label{def:b1:electronic-effects:nucleophile}
\emph{亲核试剂}\index{亲核试剂}是指向氢以外的原子
提供电子对以形成化学键的物种（\omterm{def:b1:lewis-resonance-vsepr:lewis-acid}{路易斯碱}）；
\emph{亲电试剂}\index{亲电试剂}是指接受这种电子对的物种
（\omterm{def:b1:lewis-resonance-vsepr:lewis-acid}{路易斯酸}，常为带部分正电荷的碳）。
物种的\emph{亲核性}\index{亲核性}是其作为亲核试剂的反应活性，
以\omterm{def:b1:rate-laws:order}{速率常数}衡量。
\emph{离去基团}\index{离去基团}是指当碳上的键\omterm{def:b1:electronic-effects:cleavage}{异裂}时，
带着\omterm{def:b1:lewis-resonance-vsepr:lewis}{成键电子对}一起离开的基团。
\end{definition}

\begin{proposition}[亲核性与碱性]\label{prop:b1:electronic-effects:nucleophilicity-basicity}
对于通过同一原子进攻的\omterm{def:b1:electronic-effects:nucleophile}{亲核试剂}，\omterm{def:b1:electronic-effects:nucleophile}{亲核性}
与碱性一致：\ce{CH3O-} $>$ \ce{OH-} $>$ \ce{CH3COO-} $>$ \ce{H2O}。
在\omterm{def:b1:intermolecular-forces-solvents:solvent-classes}{质子性溶剂}中，沿元素\omterm{def:b1:periodicity:table}{周期}表同一族向下，这一规律失效，
大的、易极化、\omterm{def:b1:intermolecular-forces-solvents:dissolution}{溶剂化}弱的\omterm{def:b1:electronic-effects:nucleophile}{亲核试剂}尽管碱性较弱，反应却更快
（\ce{I-} $>$ \ce{Br-} $>$ \ce{Cl-} $>$ \ce{F-}）；
对于体积庞大的碱，其\omterm{def:b1:electronic-effects:nucleophile}{亲核性}则较差。
\end{proposition}

\begin{proof}
碱性衡量的是进攻质子的平衡，
\omterm{def:b1:electronic-effects:nucleophile}{亲核性}衡量的是进攻碳的速率；两者都随电子对可用性的增加而增大，
因此同一原子的情形呈平行关系。
速率还取决于将溶剂从\omterm{def:b1:electronic-effects:nucleophile}{亲核试剂}上剥离所需的能量，
以及其电子云向拥挤碳原子的畸变程度。
小的阴离子如 \ce{F-} 与\omterm{def:b1:intermolecular-forces-solvents:solvent-classes}{质子性溶剂}分子结合得很牢，
（见\cref{ch:b1:intermolecular-forces-solvents}），
而庞大的基团阻碍其接近碳，却不妨碍其接近质子。
\end{proof}

好的\omterm{def:b1:electronic-effects:nucleophile}{离去基团}是\omterm{def:b1:predominant-reaction:strong-acid}{弱碱}，且带有电子对后仍稳定：
\ce{I-}、\ce{Br-}、\ce{Cl-}、水、磺酸根；\ce{OH-} 和 \ce{RO-} 为\omterm{def:b1:predominant-reaction:strong-acid}{强碱}，
是较差的\omterm{def:b1:electronic-effects:nucleophile}{离去基团}，这正是醇在取代前必须先被活化的原因
（见\cref{ch:b1:alcohol-activation}）。

\section{练习}

\begin{exercise}[$\star$]\label{exo:b1:electronic-effects:1}
判断 2-甲基丁烷和 2-溴-2-甲基丙烷中每个碳原子是
\omterm{def:b1:electronic-effects:carbon-class}{伯碳}、\omterm{def:b1:electronic-effects:carbon-class}{仲碳}、\omterm{def:b1:electronic-effects:carbon-class}{叔碳}还是季碳（与四个碳相连）。对于分子式为
\ce{C4H9Br} 的每种溴代物，哪一种\omterm{def:b1:electronic-effects:intermediates}{碳正离子}最容易生成？
\end{exercise}

\begin{exercise}[$\star$]\label{exo:b1:electronic-effects:2}
由 $\mathrm{p}K_a$ 值，计算氯乙酸与乙酸的
\omterm{def:b1:predominant-reaction:acidity-constant}{酸度常数}之比，以及三氟乙酸与乙酸的
\omterm{def:b1:predominant-reaction:acidity-constant}{酸度常数}之比。
\end{exercise}

\begin{exercise}[$\star$]\label{exo:b1:electronic-effects:3}
画出乙酸根离子、4-硝基苯氧根离子
（电荷位于硝基的氧原子上）以及烯丙基
正离子 \ce{CH2=CH-CH2+} 的\omterm{def:b1:lewis-resonance-vsepr:resonance}{共振结构}。
\end{exercise}

\begin{exercise}[$\star$]\label{exo:b1:electronic-effects:4}
补全图中质子转移的\omterm{def:b1:electronic-effects:curly-arrow}{弯箭头}，以及
\ce{NH3 + H3O+ -> NH4+ + H2O} 的\omterm{def:b1:electronic-effects:curly-arrow}{弯箭头}，并检查电荷。
\end{exercise}

\begin{exercise}[$\star\star$]\label{exo:b1:electronic-effects:5}
按酸性由弱到强排序：乙醇、苯酚、乙酸、
4-硝基苯酚、三氯乙酸、水（$\mathrm{p}K_a$ 为
\ce{H2O}/\ce{OH-} 电对的 14.0）。每一步都
用\omterm{def:b1:electronic-effects:inductive}{诱导效应}或\omterm{def:b1:electronic-effects:mesomeric}{共轭效应}加以论证。
\end{exercise}

\begin{exercise}[$\star\star$]\label{exo:b1:electronic-effects:6}
按碱性由弱到强排序：苯胺、甲胺、氨、吡啶。
用\omterm{def:b1:lewis-resonance-vsepr:resonance}{共振结构}解释苯胺的位置。
\end{exercise}

\begin{exercise}[$\star\star$]\label{exo:b1:electronic-effects:7}
为什么 3-硝基苯酚的酸性比 4-硝基苯酚弱，
但仍比苯酚强？画出能说明这两种说法的\omterm{def:b1:lewis-resonance-vsepr:resonance}{共振结构}。
\end{exercise}

\begin{exercise}[$\star\star$]\label{exo:b1:electronic-effects:8}
将\omterm{def:b1:electronic-effects:intermediates}{碳正离子} \ce{CH3+}、\ce{(CH3)3C+}、\ce{CH2=CH-CH2+}、
\ce{CH3CH2+}、\ce{C6H5CH2+}、\ce{CH3OCH2+} 排序，并说明理由。
\end{exercise}

\begin{exercise}[$\star\star$]\label{exo:b1:electronic-effects:9}
指出下列反应中的\omterm{def:b1:electronic-effects:nucleophile}{亲核试剂}和\omterm{def:b1:electronic-effects:nucleophile}{亲电试剂}：\ce{CH3Br + OH-}；
\ce{H2O + H+}；乙醛与氰根离子；\ce{BF3 + NH3}。画出
\omterm{def:b1:electronic-effects:curly-arrow}{弯箭头}。
\end{exercise}

\begin{exercise}[$\star\star\star$]\label{exo:b1:electronic-effects:10}
某溶液含有 \qty{0.010}{mol/L} 的三氯乙酸。判断
\omterm{def:b1:predominant-reaction:strong-acid}{弱酸}公式 $\mathrm{pH} = \frac12(\mathrm{p}K_a - \log C)$ 是否成立？
准确计算该溶液的 pH 和\omterm{def:b1:predominant-reaction:dissociation}{解离度}。
\end{exercise}

\begin{exercise}[$\star\star\star$]\label{exo:b1:electronic-effects:11}
在水中测得的醇的 $\mathrm{p}K_a$ 由甲醇
（15.5）到乙醇（15.9）、2-丙醇（17.1）以及 2-甲基-2-丙醇
（19.2）递增。给出两种解释：一种通过烷基对醇盐的\omterm{def:b1:electronic-effects:inductive}{诱导效应}，
另一种通过醇盐的\omterm{def:b1:intermolecular-forces-solvents:dissolution}{溶剂化}。
% ledger: pka:methanol, pka:ethanol, pka:2-propanol, pka:tBuOH
\end{exercise}

\begin{exercise}[$\star\star\star$]\label{exo:b1:electronic-effects:12}
解释为什么酰胺既不在氮原子上显碱性，也不具有\omterm{def:b1:electronic-effects:nucleophile}{亲核性}，
而胺两者兼有，请利用酰胺的\omterm{def:b1:lewis-resonance-vsepr:resonance}{共振结构}加以说明。
在\omterm{def:b1:predominant-reaction:strong-acid}{强酸}中，酰胺的哪个原子会被质子化？
\end{exercise}

\section{周末问题：三氟乙酸为何如此之强}

\begin{problem}[{周末问题——卤代酸的诱导效应系列、
羧酸根与醇盐的比较、三种硝基苯酚、
优势反应的计算，以及三氟乙酸与乙酸
的酸度常数之比}]
\label{pb:b1:electronic-effects:1}
数据（\qty{25}{\celsius} 下的 $\mathrm{p}K_a$）：乙酸 4.76，
氯乙酸 2.87，二氯乙酸 1.35，三氯乙酸 0.51，
三氟乙酸 0.52，乙醇 15.9，苯酚 9.99，2-硝基苯酚 7.23，
3-硝基苯酚 8.36，4-硝基苯酚 7.16；$\mathrm{p}K_e = 14.00$。

\medskip
\noindent\textbf{第一部分——\omterm{def:b1:electronic-effects:inductive}{诱导效应}系列。}
\begin{enumerate}
  \item 计算从乙酸到三氯乙酸的每一步中
        $K_a$ 之比。
  \item 每个氯的作用是否相同？请加以说明。
  \item 用\omterm{def:b1:electronic-effects:inductive}{诱导效应}解释这一变化趋势。
  \item 为什么较长碳链上离 COOH 最远的碳原子所连的氯
        几乎没有影响？
  \item 氟的\omterm{def:b1:periodicity:electronegativity}{电负性}比氯强。三氟乙酸与三氯乙酸的比较
        说明了什么，以及为什么
        难以精确测量这类 $\mathrm{p}K_a$？
  \item 在 pH 为 3 时，这些酸各以何种形式存在？
\end{enumerate}

\medskip
\noindent\textbf{第二部分——羧酸根与醇盐。}
\begin{enumerate}[resume]
  \item 画出乙酸根离子的两种\omterm{def:b1:lewis-resonance-vsepr:resonance}{共振结构}。
  \item 它们对该离子中两个 C--O 键的键长给出了怎样的预测？
  \item 为什么乙氧根离子没有等价的\omterm{def:b1:lewis-resonance-vsepr:resonance}{共振结构}？
  \item 计算乙酸与乙醇的 $K_a$ 之比。
  \item 将苯酚排在二者之间，并用苯氧根离子的
        \omterm{def:b1:lewis-resonance-vsepr:resonance}{共振结构}加以说明。
  \item 乙醇、苯酚和乙酸中，哪些能与碳酸氢盐溶液几乎
        完全反应（$\mathrm{p}K_a$ 为 6.37 的
        \ce{CO2,H2O}/\ce{HCO3-} 电对）？
\end{enumerate}

\medskip
\noindent\textbf{第三部分——硝基苯酚。}
\begin{enumerate}[resume]
  \item 按酸性对三种硝基苯酚和苯酚排序。
  \item 画出 4-硝基苯氧根离子的一种\omterm{def:b1:lewis-resonance-vsepr:resonance}{共振结构}，其中
        电荷位于硝基上。
  \item 说明 3-硝基苯氧根离子不存在这样的\omterm{def:b1:lewis-resonance-vsepr:resonance}{共振结构}。
  \item 为什么 3-硝基苯酚的酸性仍强于苯酚？
  \item 2-硝基苯酚的酸性比 4-硝基苯酚略弱，尽管两者都
        具有\omterm{def:b1:electronic-effects:mesomeric}{共轭效应}。请提出一种解释，涉及其
        羟基与邻近硝基的作用。
  \item 在 pH 为 7.5 时，哪些硝基苯酚主要以阴离子形式存在，
        这为何使 4-硝基苯酚可作\omterm{def:b1:titration-methods:indicator}{显色指示剂}？
\end{enumerate}

\medskip
\noindent\textbf{第四部分——水中的三氟乙酸。}
\begin{enumerate}[resume]
  \item 计算浓度为 \qty{0.010}{mol/L} 的三氟乙酸溶液的 pH
        （不要假设为弱解离）。
  \item 计算浓度为 \qty{0.010}{mol/L} 的乙酸溶液的 pH。
  \item 写出三氟乙酸与乙酸钠的反应式，
        并计算其平衡常数。
  \item 给出比值 $K_a(\ce{CF3COOH})/K_a(\ce{CH3COOH})$，分别以
        10 的幂次和数值表示。
\end{enumerate}
\end{problem}
